\section{Das A und O eines Gedichtes}
\stepcounter{ctrpoeminyear}
%\addcontentsline{toc}{section}{\thechapter.\arabic{ctrpoeminyear} \; Das A und O eines Gedichtes}
%addcontentsline is necessary for the correct page number to appear in the toc




\begin{strophe}
\vers{Das Alphabet mit A beginnt.}
\vers{Mit einundzwanzig Konsonanten geschmückt,}
\vers{doch nur fünf Vokale, ich werde verrückt!}
\vers{Das lernt man früh, wie jedes Kind.}
\end{strophe}

\begin{strophe}
\vers{Die Fünf die jeder zu kennen vermag}
\vers{werden jetzt genannt, Schlag auf Schlag.}
\end{strophe}

\begin{strophe}
\vers{Kommt einem ein Geistesblitz,}
\vers{so ruft man: \ttquotl Ah!\ttquotr}
\vers{Das ist doch klar.}
\vers{Dann springt man auf von seinem Sitz.}
\end{strophe}

\begin{strophe}
\vers{Den nächsten Vokal findet man selten allein,}
\vers{doch spielst du gern im Schnee,}
\vers{antwortest du auf: \ttquotl Ist dir warm?\ttquotr}
\vers{Mit 'nem einfachen \ttquotl Nee\ttquotr.}
\vers{Natürlich mit etwas Charme.}
\vers{Das sagt man schon wenn man ist klein.}
\end{strophe}

\begin{strophe}
\vers{Ist etwas Ekelhaft,}
\vers{das hört man fast nie,}
\vers{doch dann ertönt es: \ttquotl Ihh!\ttquotr}
\vers{Das schreit man mit aller Kraft.}
\end{strophe}

\begin{strophe}
\vers{Schmerzt dir der Bauch}
\vers{Kommt aus deinem Mund ein leises \ttquotl Oh\ttquotr .}
\vers{Ein Tipp von mir: Versuch's mal auf dem Klo,}
\vers{aber bitte setz' dich drauf.}
\end{strophe}

\begin{strophe}
\vers{Nun kommen wir zur Bewunderung:}
\vers{Mit einem \ttquotl Uh\ttquotr\ staunst du nicht schlecht,}
\vers{doch, Ja, dieses Gedicht ist echt.}
\vers{Hör' ich da vielleicht ein \ttquotl Uh\ttquotr\ aus Verwunderung?}
\end{strophe}

14.05.2018 (Revision 2, bzw. Version 3)

%Christoph Koloska




